\bookchapter{Κεφάλαιο 17: Σύνδεση κινητού $\mu$ε άλλες συσκευές}

\chapterintro{Το κινητό σας είναι σούπερ τηλεχειριστήριο και κέντρο δεδο$\mu$ένων. Δεν κάνει $\mu$όνο κλήσεις, συνδέει υπολογιστή, τηλεόραση, ρολόι, ηχείο, αυτοκίνητο, ώστε να $\mu$ην είναι νησιά. Θα δού$\mu$ε πώς να τα κάνου$\mu$ε να “συνεργάζονται”.}

\booksection{I. Συνεργασία $\mu$ε υπολογιστή: κινητό και υπολογιστής δεν είναι πια δύο κόσ$\mu$οι}

Ειλικρινά, παλιά ενοχλούσε να βάλω φωτογραφία σε υπολογιστή, χρειαζόταν καλώδιο και ανα$\mu$ονή. Αργότερα κατάλαβα: δεν χρειάζεται, η καλή διασύνδεση είναι 10 φ faster.

\begin{itemize}
  \item \textbf{Apple οικοσύστη$\mu$α (AirDrop + universal clipboard)}: iPhone+Mac: AirDrop. Share AirDrop Mac, φωτογραφία α$\mu$έσως, χωρίς data/Wi-Fi, BT+Wi-Fi Direct. "Universal clipboard": copy στο κινητό, paste στο Mac, $\mu$αγεία.
  \item \textbf{Android+Windows (Σύνδεση $\mu$ε Windows)}: Microsoft "Link to Windows", Samsung built-in, άλλα από store. Ειδοποιήσεις, SMS, φωτογραφίες, κλήσεις σε PC. 2 χρόνια χρήσης, χωρίς να σκύβω.
  \item \textbf{Συνεργασία πολλαπλών οθονών επωνυ$\mu$ίας}: Huawei "Multi-screen" και Xiaomi "Mijia":προβολή οθόνη σε υπολογιστή, drag-drop αρχεία, Douyin, WeChat. Copy-paste, smooth.
\end{itemize}

\booksection{II. Cast στην τηλεόραση: όλη η οικογένεια}

Πρωτοχρονιά, η ξαδέρφη ήθελε να βάλει βίντεο ταξιδιού στην τηλεόραση. Δοκι$\mu$άσα$\mu$ε και τελικά ήταν απλό:

\begin{itemize}
  \item \textbf{iPhone τηλεόραση}: Σύρετε δεξιά πάντω, εικονίδιο δύο κουτιά (Screen Mirroring), βρείτε τηλεόραση. Ίδιο Wi-Fi και AirPlay (οι περισσότερες smart TVs).
  \item \textbf{Android Τηλεόραση}: Διαφορετικά ονό$\mu$ατα: Huawei "wireless cast", Xiaomi "cast", OPPO/vivo "multi-screen". Στο quick panel ή ψάξτε "cast".
  \item \textbf{Cast απευθείας από App βίντεο}: Αν θέλετε iQiyi/B σταθ$\mu$ός/YouTube σε TV, απλό — εικονίδιο rectangle+Wi-Fi, πατήστε, επιλέξτε TV. Το κινητό ελεύθερο.
  \item \textbf{Τι κάνετε αν αποτύχει το cast}: 90\% δεν είναι στο ίδιο Wi-Fi. 10\% η τηλεόραση δεν έχει ανοιχτό το cast, ψάξτε "multi-screen" ή "AirPlay".
\end{itemize}

\booksection{III. Με έξυπνο ρολόι/bracelet: διαχειριστής υγείας στον καρπό}

Αγόρασα στη $\mu$α$\mu$ά $\mu$ου ένα βραχιόλι δραστηριότητας, στην αρχή είπε “αφού δεν γυ$\mu$νάζο$\mu$αι, γιατί να το αγοράσω;”. Μετά από δύο εβδο$\mu$άδες χρήσης, ανακάλυψε ότι κοι$\mu$όταν $\mu$όνο λίγο πάνω από πέντε ώρες την η$\mu$έρα και άρχισε να προσαρ$\mu$όζει ενεργά το πρόγρα$\mu$$\mu$ά της. Αυτή είναι η πιο πρακτική αξία των ρολογιών/βραχιολιών — $\mu$ετατρέπουν τα αόρατα δεδο$\mu$ένα υγείας σε ορατούς αριθ$\mu$ούς.

\begin{itemize}
  \item \textbf{Φίλτρο ειδοποιήσεων}: Σε τσέπη, σηκώστε καρπό — boss ή Pinduododo promo, $\mu$ια κίνηση.
  \item \textbf{Συγχρονισ$\mu$ός δεδο$\mu$ένων υγείας}. Ρολόι 24/7: καρδιακός, SpO2, ύπνος, βή$\mu$ατα, στο health App. Τάσεις.
  \item \textbf{Συγχρονισ$\mu$ός κάρτας $\mu$εταφοράς/ελέγχου πρόσβασης}: Κάρτες στο κινητό, sync σε ρολόι. Στο $\mu$ετρό χωρίς κινητό, καρπός στο turnstile, "$\mu$έλλον".
  \item \textbf{Εργαλείο εύρεσης κινητού}. Ρολόι "βρες κινητό", πατήστε, χτυπάει. Συχνή χρήση.
\end{itemize}

\booksection{IV. Με το έξυπνο σπίτι: το τηλέφωνο ως ο “κεντρικός διακόπτης” του σπιτιού}

Το smart home $\mu$οιάζει $\mu$υστήριο, αλλά είναι τρία πράγ$\mu$ατα: τηλεχειρισ$\mu$ός, αυτο$\mu$ατισ$\mu$οί, φωνητικός έλεγχος.

Μερικές δικές $\mu$ου χρήσεις, να δείτε πόσο πρακτικό είναι:

\begin{itemize}
  \item \textbf{Κλι$\mu$ατισ$\mu$ός}: Καλοκαίρι, ανοίξτε AC από το κινητό, δροσερά. Με AC mate πρίζα 100+, χωρίς νέο AC.
  \item \textbf{Φως}. Πριν ύπνο, πείτε "Hey Siri/Xiaoai, σβήσε φως".
  \item \textbf{Ρο$\mu$πότ σκουπίσ$\mu$ατος}: Πριν βγείτε, πατήστε "έναρξη σκούπισ$\mu$α", γυρνάτε καθαρό. Και καθη$\mu$ερινό timer.
  \item \textbf{Κά$\mu$ερα}: Με ηλικιω$\mu$ένους/παιδιά/κατοικίδια, δείτε σπίτι από παντού.
\end{itemize}

Ποιο οικοσύστη$\mu$α; Με $\mu$ια φράση:

\begin{itemize}
  \item Με iPhone προτι$\mu$ήστε Apple HomeKit (πιο ακριβά αλλά καλύτερη ιδιωτικότητα).
  \item Πολλές συσκευές, value Xiaomi Mijia πιο πλήρες, φώτα, πρίζες, ρο$\mu$πότ σκούπες.
  \item Με Huawei προτι$\mu$ήστε HarmonyOS, καλύτερη ε$\mu$πειρία $\mu$ε Huawei.
\end{itemize}

\booksection{V. Μπλουτούθ ακουστικά και ηχεία: η ασύρ$\mu$ατη ελευθερία}

Η αλλαγή από ενσύρ$\mu$ατα σε $\mu$πλουτούθ ήταν η πιο αισθητή αναβάθ$\mu$ιση. Δεν βγάζω πια κουβάρι, ούτε φοβά$\mu$αι ότι το καλώδιο θα πιαστεί στη λαβή.

\begin{itemize}
  \item \textbf{Η σύζευξη είναι απλή}. Πρώτη φορά pairing mode, BT search, ένα πάτη$\mu$α. Μετά αυτό$\mu$ατο.
  \item \textbf{Ίδια επωνυ$\mu$ία έχει πλεονέκτη$\mu$α}: AirPods $\mu$ε iPhone δείχνουν popup $\mu$παταρίας και αλλάζουν συσκευές. FreeBuds $\mu$ε Huawei, ακουστικά Xiaomi έχουν παρό$\mu$οια.
  \item \textbf{Μπλουτούθ ηχείο}: Σε party/camping, $\mu$πλουτούθ ηχείο διπλασιάζει τη διάθεση. Το κινητό DJ.
  \item \textbf{Προσοχή στην καθυστέρηση}. Παιχνίδι $\mu$ε $\mu$πλουτούθ = καθυστέρηση (η $\mu$ισή δευτ.). Ευαίσθητοι: ενσύρ$\mu$ατα ή "low latency" mode.
\end{itemize}

\begin{tcbinfo}{Τεχνολογία UWB — πιο ακριβής «χωρική αίσθηση» από το Bluetooth}
Bluetooth, UWB(Ultra-Wideband,).\textbf{Ακριβής $\mu$έτρηση απόστασης και χωρικός εντοπισ$\mu$ός}——Συσκευή, .Bluetooth“”,UWB“1.2”.Εφαρ$\mu$ογή:

\begin{itemize}
  \item \textbf{AirTag/εύρεση $\mu$ε ακρίβεια}:iPhone εύρεση AirTag ε$\mu$φάνιση "7,"Πλοήγηση, UWB.Bluetooth"".
  \item \textbf{Ψηφιακό κλειδί αυτοκινήτου}:UWBiPhone ναυαρχίδα, αυτό$\mu$ατο Ξεκλείδω$\mu$α, αυτό$\mu$ατο, .
  \item \textbf{Εσωτερική πλοήγηση}:,UWB""Πλοήγηση, GPS.
  \item \textbf{Υποστηριζό$\mu$ενα $\mu$οντέλα}:iPhone 11UWB(U1/U2); vivo, OPPO ναυαρχίδα.
\end{itemize}

\end{tcbinfo}

\begin{tcbinfo}{LE Audio και Auracast — η νέα γενιά ήχου Bluetooth}
Το κλασικό Bluetooth audio (A2DP) δύο δεκαετιών αντικαθίσταται\textbf{LE Audio}(χα$\mu$ηλής κατανάλωσης ήχος), η νέα γενιά από Bluetooth 5.2:

\begin{itemize}
  \item \textbf{Λιγότερη ενέργεια, καλύτερος ήχος}:LC3 Κωδικοποίηση, SBC, Ακουστικά.
  \item \textbf{Διπλό κανάλι ήχου}: Ανεξάρτητη $\mu$ετάδοση ανά ακουστικό, εναλλαγή χωρίς αποσύνδεση.
  \item \textbf{Auracast εκπο$\mu$πή ήχου}: Η λειτουργία-ση$\mu$αία του LE Audio. Ένα κινητό/TV $\mu$εταδίδει ταυτόχρονα σε\textbf{άπειρα}Ακουστικά Ήχος——, Ακουστικά.
  \item \textbf{Τρέχουσα κατάσταση}:2024-2026 ναυαρχίδα Ακουστικά LE Audio, Συσκευή αναβάθ$\mu$ιση Ξεκλείδω$\mu$α.Ακουστικά LE Audio/Auracast.
\end{itemize}

\end{tcbinfo}

\booksection{VI. Με σύστη$\mu$α αυτοκινήτου: ασφαλής συνεργάτης}

Το να κοιτάτε το κινητό ενώ οδηγείτε είναι παράνο$\mu$ο και θανατηφόρο. Αλλά $\mu$ερικές φορές χρειάζεστε πλοήγηση, απάντηση κλήσεων, αναπαραγωγή $\mu$ουσικής. Τότε το CarPlay (Apple) και το Android Auto (Android) είναι χρήσι$\mu$α.

\begin{itemize}
  \item \textbf{Πώς συνδέεστε}: Αν το αυτοκίνητο υποστηρίζει, καλώδιο σε USB, οθόνη αλλάζει σε CarPlay/Android Auto. Κάποια ασύρ$\mu$ατα.
  \item \textbf{Τι $\mu$πορεί να κάνει}: Navigation (καλύτερο από auto), κλήσεις (τι$\mu$όνι), $\mu$ουσική, φωνητικό WeChat (Siri/Xiao Ai).
  \item \textbf{Λειτουργεί και σε παλαιά αυτοκίνητα}: Αν παλιό, $\mu$ερικές εκατοντάδες για external CarPlay οθόνη, παρό$\mu$οια.
\end{itemize}

\booksection{VII. Ασφάλεια και ιδιωτικότητα στη διασύνδεση}

Όσο περισσότερες συσκευές, τόσο $\mu$εγαλύτερη επιφάνεια ρίσκου. Ση$\mu$εία:

\begin{itemize}
  \item Πριν κάνετε cast σε δη$\mu$όσιο χώρο, ελέγξτε τις ειδοποιήσεις για προσωπικά $\mu$ηνύ$\mu$ατα.
  \item Πριν δώσετε ή πουλήσετε συσκευή, αποσυνδέστε τη — στις ρυθ$\mu$ίσεις βρείτε τη συσκευή και επιλέξτε "Αφαίρεση " ή "Ακύρωση σύζευξης".
  \item Αγοράστε κά$\mu$ερες smart home $\mu$ε $\mu$άρκα και κρυπτογράφηση, όχι φθηνές. Πάρα πολλές ειδήσεις για κά$\mu$ερες που έγιναν live από χάκερ.
\end{itemize}

\begin{tcbpractice}{Δοκι$\mu$ή: κάντε cast στην τηλεόραση και στήστε $\mu$ια έξυπνη σκηνή}
\begin{enumerate}
  \item Σιγουρευτείτε ότι κινητό και τηλεόραση είναι στο ίδιο Wi-Fi.
  \item Βρείτε τη λειτουργία "Screen Mirroring" ή "AirPlay" στην τηλεόραση και ενεργοποιήστε την.
  \item Στο κινητό: iPhone — από το Κέντρο Ελέγχου πατήστε «Κατοπτρισ$\mu$ός οθόνης». Android — σύρετε κάτω και βρείτε «Ασύρ$\mu$ατη προβολή».
  \item Κάντε cast την αγαπη$\mu$ένη σας φωτογραφία στην τηλεόραση.
  \item Εάν έχετε έξυπνες οικιακές συσκευές, δοκι$\mu$άστε να δη$\mu$ιουργήσετε $\mu$ια "Λειτουργία εκτός σπιτιού" στην εφαρ$\mu$ογή: ένα κλικ για να σβήσουν τα φώτα, να κλείσει το κλι$\mu$ατιστικό, να ξεκινήσει η ρο$\mu$ποτική σκούπα.
\end{enumerate}

\end{tcbpractice}

\begin{tcbinfo}{Περίπτωση 1: Πρακτική βελτιστοποίηση αυτονο$\mu$ίας}
\begin{itemize}
  \item \textbf{Σενάριο}: Ο χρήστης του Xiaomi 14 έχει κακή αυτονο$\mu$ία, εξαντλείται σε $\mu$ισή $\mu$έρα.
  \item \textbf{Βή$\mu$ατα}:(1) Ρυθ$\mu$ίσεις Μπαταρία Προβολή Κατανάλωση ενέργειας (2) Παρασκήνιο App (3) απενεργοποίηση 5G 4G (4) περιορισ$\mu$ός Παρασκήνιο (5) ενεργοποίηση Εξοικονό$\mu$ηση ενέργειας (6) Φόρτιση.
  \item \textbf{Ανα$\mu$ενό$\mu$ενο αποτέλεσ$\mu$α}: Αυτονο$\mu$ία +30-50\%, ευανάγνωστα δεδο$\mu$ένα υγείας $\mu$παταρίας.
\end{itemize}

\end{tcbinfo}
